% =====================================================================
\chapter{Implementation and Integration}
\label{ch:implementation}
% =====================================================================

This chapter describes how the methodology of Chapter~\ref{ch:methodology} is
realised as a research framework: the backbone and software stack
(\S\ref{sec:impl-framework}), the document-IE data pipeline
(\S\ref{sec:impl-pipeline}), and the system architecture, including the
per-component instrumentation and experiment tracking
(\S\ref{sec:impl-architecture}).

\section{Framework Selection}
\label{sec:impl-framework}

\paragraph{Backbone.} The primary backbone is LayoutLMv3-base
(\texttt{microsoft/layoutlmv3-base}, $\approx$\,133M parameters, 12 layers, hidden
size 768). It was chosen because it is document-specific, encoder-only, and small
enough to run the full experimental grid---ten baselines and the selected method
across three scenarios and three seeds---within the compute budget, while exposing
the four separable components central to the diagnosis. LiLT-base and BROS-base
serve as secondary backbones for a generalisation study on a reduced method
subset. Table~\ref{tab:backbone-comparison} compares the candidates.

\begin{table}[t]
  \centering
  \caption[Backbone comparison]{Backbones considered for the study. LayoutLMv3-base
  is the primary subject; LiLT and BROS test generalisation of the diagnosis; large
  multimodal LLMs are noted as comparison targets but are out of scope for the
  controlled per-component study.}
  \label{tab:backbone-comparison}
  \begin{tabular}{llll}
    \toprule
    \textbf{Backbone} & \textbf{Type} & \textbf{Scale} & \textbf{Role} \\
    \midrule
    LayoutLMv3-base & encoder-only, tri-modal & $\approx$\,133M & primary \\
    LiLT-base       & encoder-only, decoupled layout & $\approx$\,130M & secondary \\
    BROS-base       & encoder-only, spatial & $\approx$\,110M & secondary \\
    InternVL3-8B    & multimodal LLM & 8B & comparison target (out of scope) \\
    \bottomrule
  \end{tabular}
\end{table}

\paragraph{Software stack.} The framework is implemented in Python with PyTorch
and the HuggingFace Transformers library for the backbones, \texttt{seqeval} for
BIO-aware entity-F1, and Hydra for composable configuration. Experiments are
launched through a single entry point and configured by composition, e.g.
\begin{quote}\small
\texttt{python scripts/train.py method=ewc scenario=cil\_cord seed=42
method.lambda\_=1000}
\end{quote}
so that every run is fully specified by its config and reproducible.

\section{Document-IE Pipeline Integration}
\label{sec:impl-pipeline}

Each document is loaded as OCR tokens with normalised bounding boxes
($[0,1000]^4$) and a $224\times224$ page image. Tokens are WordPiece-tokenised and
aligned to their boxes; sub-tokens inherit their word's box and only the first
sub-token carries the label, with the remainder masked to the ignore index. Image
patches are extracted by the backbone's vision module. The same preprocessing is
shared across all datasets so that the only thing that changes between continual
tasks is the data and label space, not the pipeline. For the diagnostic conditions
(C1--C4) a modality mask is applied at the input layer---zeroing image pixels,
replacing text with padding tokens, or zeroing layout positions---through the
backbone wrapper's \texttt{modality\_mask} forward argument rather than a separate
ablated model, so the four conditions differ only in which input pathways are
active.

\paragraph{Data integrity and validation.} Three preprocessing invariants are
enforced and checked so that the forgetting signal is not confounded by data
artefacts. First, WordPiece token--box alignment is validated at load time: every
sub-token is tied to exactly one word box, the first sub-token of a word carries
that word's BIO label, and the trailing sub-tokens are set to the $-100$ ignore
index so that the loss and the entity-F1 metric see one prediction per word.
Second, bounding boxes are normalised to the $[0,1000]^4$ integer range expected by
the backbone's 2D position embeddings, with clamping to guard against
out-of-range or degenerate (zero-area) boxes. Third, and most importantly for a
controlled diagnosis, the pipeline consumes the dataset-provided OCR transcriptions
and boxes directly and never re-runs an OCR engine; this removes OCR quality as a
confounder, so that any measured drift is attributable to continual training of the
encoder rather than to changing recognition behaviour across tasks.

\section{System Architecture Overview}
\label{sec:impl-architecture}

\begin{figure}[t]
  \centering
  \includegraphics[width=0.95\linewidth]{system_architecture.pdf}
  \caption[System architecture]{System architecture. The continual-learning core
  exposes a uniform method interface; the backbone wrapper exposes extension points
  for classifier expansion, prompts, and LoRA banks; the instrumentation layer
  records per-layer CKA and per-group Fisher; all runs report to the experiment
  tracker. Data flows from the shared document pipeline through the wrapped backbone
  to the method-specific heads, while the instrumentation layer taps the backbone at
  task boundaries without altering the training path.}
  \label{fig:system-architecture}
\end{figure}

Figure~\ref{fig:system-architecture} gives the overall picture: a thin
continual-learning core orchestrates the per-task loop, the backbone wrapper hides
architecture-specific details behind a stable set of extension points, the document
pipeline feeds both, and a decoupled instrumentation layer observes the backbone
without perturbing it. The remainder of this section describes each block in turn.

\subsection{Continual-Learning Core}
\label{sec:impl-cl-core}

All methods implement a common \texttt{ContinualMethod} interface with four hooks:
\texttt{before\_task} (e.g.\ expand the classifier head or allocate a LoRA bank),
\texttt{train\_task} (the training loop), \texttt{after\_task} (update the replay
buffer, compute Fisher, freeze weights), and \texttt{evaluate} (evaluate on all
tasks seen so far). This interface lets regularisation, replay, prompt-based, and
LoRA-based methods share the same training and evaluation harness, which is
essential for like-for-like comparison.

\subsection{Per-Component Instrumentation}
\label{sec:impl-instrumentation}

The diagnostic of \S\ref{sec:method-diagnostic} is implemented as a set of hooks
that, at each task boundary, capture activations at the tracked layers (for CKA)
and accumulate squared gradients per parameter group (for Fisher), writing
per-run JSON dumps that are later aggregated into a long-format dataframe for
analysis and plotting. The instrumentation is decoupled from the methods so that
it can be attached to any run, including the naive-sequential pilot runs. Four
mechanisms make up the layer:

\begin{itemize}
  \item \textbf{Representational drift (CKA).} Forward hooks capture per-layer
  activations on a \emph{fixed} probe set---the same held-out documents are used at
  every task boundary so that representations are compared on identical inputs.
  Linear CKA between the activations of two checkpoints scores how much each layer's
  representation has changed; because it is invariant to orthogonal transforms and
  isotropic scaling, a drop signals genuine representational reorganisation rather
  than a trivial rotation. The probe activations are gathered through a hook helper
  and scored layer by layer.
  \item \textbf{Parameter importance (per-group Fisher).} Over a task's data the
  diagonal empirical Fisher accumulates squared gradients of the log-likelihood with
  respect to each parameter, summed within the eight named parameter groups
  (text word embeddings, 2D layout-position embeddings, image patch embedding,
  text attention, visual attention, cross-modal fusion, feed-forward, and
  classifier). Comparing the per-group Fisher before and after a task gives a
  per-component importance-drop profile that complements the CKA view.
  \item \textbf{Classifier-head expansion.} When a new task introduces new BIO
  labels, the classification head is grown to cover them while the rows
  corresponding to previously seen labels are preserved, so that old-label logits
  are not reset by the arrival of new classes. Expansion is performed once by the
  training loop before each task, and the label$\leftrightarrow$index maps are
  updated in lockstep.
  \item \textbf{Modality ablation.} The pilot conditions are realised by zeroing
  input pathways \emph{in place} rather than by instantiating separate models: the
  text, visual, and layout streams are masked individually through a single forward
  argument, so C1--C4 share identical weights and differ only in which modalities
  reach the encoder. This is what makes the per-condition comparison
  apples-to-apples.
\end{itemize}

\paragraph{Code pointers.} For reproducibility, the components above map directly
onto the source: linear CKA and the activation hooks live in
\texttt{doccl/eval/cka.py}; the diagonal and per-group Fisher in
\texttt{doccl/eval/fisher.py}; classifier expansion, the eight-group
\texttt{param\_groups} mapping, and the \texttt{modality\_mask} forward argument in
\texttt{doccl/models/layoutlm\_wrapper.py}; and the method lifecycle hooks in
\texttt{doccl/methods/base.py}. The single Hydra entry point
\texttt{scripts/train.py} wires these together and resolves each method from its
registry.

\subsection{Datasets and Scenario Builders}
\label{sec:impl-datasets}

Dataset loaders expose FUNSD, CORD, and SROIE in the shared format, and scenario
builders assemble them into continual sequences (class-incremental, domain-incremental,
mixed) through a registry, so that a scenario is selected by name from the config.

\subsection{Compute and Deployment}
\label{sec:impl-compute}

Experiments run on a single RTX 4090 (cloud), with a local RTX 2060 reserved for
debugging only; paper-result experiments are never run on the debug machine.
Determinism is enforced (fixed seeds, deterministic cuDNN) so that CKA and Fisher
patterns reflect signal rather than run-to-run noise.

\subsection{Experiment Tracking and Monitoring}
\label{sec:impl-tracking}

All runs report to Weights \& Biases under a single project, with a structured tag
schema (\texttt{\{phase\}/\{scenario\}/\{method\}/seed\{seed\}}; pilot runs use
\texttt{pilot/\{condition\}/seed\{seed\}}), so that the full grid is queryable and
the result tables of Chapter~\ref{ch:results} can be regenerated from the tracked
logs.

\section{Summary}
\label{sec:impl-summary}

The framework pairs a small, component-separable backbone with a uniform
continual-method interface, a shared document pipeline, decoupled per-component
instrumentation, and structured experiment tracking. This design makes the
diagnosis attachable to any run and the baseline comparison strictly like-for-like,
which the next chapter relies upon.
